\section{12. Líneas de evolución}
La siguiente etapa debe consolidar la operación existente y ampliar progresivamente la capacidad de EBP sin comprometer la estabilidad del servicio.

\begin{tabularx}{\textwidth}{|L{0.30\textwidth}|Y|}\hline
\rowcolor{Navy}\THead{Línea} & \THead{Alcance} \\ \hline
Consolidación operacional & Mayor validación física, historial, recuperación y seguimiento de eventos. \\ \hline
\rowgray Análisis audiovisual & Detección automática de ausencia de video/audio, congelamiento, pantalla negra y otras condiciones de contenido. \\ \hline
Correlación de incidencias & Uso conjunto de evidencia de señal, sesión, protocolo, infraestructura y red para ayudar a localizar el origen de una falla. \\ \hline
\rowgray Multiviewer & Visualización simultánea y en tiempo real de las señales entrantes para supervisión operacional. \\ \hline
Nuevas fuentes y Readers & Incorporación controlada de nuevas señales y destinos. \\ \hline
\rowgray Multi-nodo & Distribución geográfica de funciones y aumento de capacidad/resiliencia. \\ \hline
Transcodificación dedicada & Conversión de formatos y bitrates mediante infraestructura separada. \\ \hline
\rowgray Servicios a terceros & Evaluación futura de distribución para otros proveedores, sujeta a capacidad, seguridad, acuerdos comerciales y derechos de contenido. \\ \hline
\end{tabularx}

El proyecto se encuentra en fase de consolidación y validación operacional. Antes de una expansión comercial deben dimensionarse aspectos como concurrencia, redundancia, seguridad, acuerdos de nivel de servicio, capacidad de infraestructura y derechos de distribución.
